% Template Beamer File for Differential Equations Chapter
% Copy this file and modify the content for your specific chapter

%% Preamble - Common setup for all beamer presentations
%% Include this file in your chapter presentations

\documentclass[aspectratio=169, 12pt]{beamer}

% Packages
\usepackage[utf8]{inputenc}
\usepackage[T1]{fontenc}
\usepackage{amsmath}
\usepackage{amssymb}
\usepackage{amsthm}
\usepackage{mathtools}
\usepackage{graphicx}
\usepackage{xcolor}
\usepackage{tcolorbox}
\usepackage{tikz}
\usetikzlibrary{shadows, arrows, shapes}
\usepackage{hyperref}
\usepackage{listings}
\usepackage{geometry}

% Custom colors (import from custom_colors.tex if available)
\definecolor{defcolor}{RGB}{230, 245, 250}
\definecolor{defborder}{RGB}{0, 102, 204}
\definecolor{thmcolor}{RGB}{245, 230, 250}
\definecolor{thmborder}{RGB}{153, 51, 153}
\definecolor{excolor}{RGB}{240, 255, 240}
\definecolor{exborder}{RGB}{34, 139, 34}
\definecolor{probcolor}{RGB}{255, 250, 240}
\definecolor{probborder}{RGB}{255, 140, 0}

% Beamer Theme
\usetheme{Madrid}
\usecolortheme{default}
\setbeamertemplate{navigation symbols}{}
\setbeamertemplate{footline}[frame number]

% tcolorbox environment definitions

%% Definition Box
\newtcolorbox{definitionbox}[1]{
  colback=defcolor,
  colframe=defborder,
  coltitle=white,
  colbacktitle=defborder,
  title=#1,
  fonttitle=\bfseries\large,
  boxrule=2pt,
  arc=5pt,
  left=8pt,
  right=8pt,
  top=8pt,
  bottom=8pt,
  shadow={2pt}{-2pt}{0pt}{black!20},
  enhanced
}

%% Theorem Box
\newtcolorbox{theorembox}[1]{
  colback=thmcolor,
  colframe=thmborder,
  coltitle=white,
  colbacktitle=thmborder,
  title=#1,
  fonttitle=\bfseries\large,
  boxrule=2pt,
  arc=5pt,
  left=8pt,
  right=8pt,
  top=8pt,
  bottom=8pt,
  shadow={2pt}{-2pt}{0pt}{black!20},
  enhanced
}

%% Example Box
\newtcolorbox{examplebox}[1]{
  colback=excolor,
  colframe=exborder,
  coltitle=white,
  colbacktitle=exborder,
  title=#1,
  fonttitle=\bfseries\large,
  boxrule=2pt,
  arc=5pt,
  left=8pt,
  right=8pt,
  top=8pt,
  bottom=8pt,
  shadow={2pt}{-2pt}{0pt}{black!20},
  enhanced
}

%% Problem Box
\newtcolorbox{problembox}[1]{
  colback=probcolor,
  colframe=probborder,
  coltitle=white,
  colbacktitle=probborder,
  title=#1,
  fonttitle=\bfseries\large,
  boxrule=2pt,
  arc=5pt,
  left=8pt,
  right=8pt,
  top=8pt,
  bottom=8pt,
  shadow={2pt}{-2pt}{0pt}{black!20},
  enhanced
}

% Custom commands
\newcommand{\RR}{\mathbb{R}}
\newcommand{\CC}{\mathbb{C}}
\newcommand{\NN}{\mathbb{N}}
\newcommand{\ZZ}{\mathbb{Z}}
\newcommand{\dd}[1]{\frac{d}{d#1}}
\newcommand{\dydx}{\frac{dy}{dx}}
\newcommand{\dd2ydx2}{\frac{d^2y}{dx^2}}

% Title page formatting
\title[Your Chapter Title]{Your Chapter Title}
\author{Author Name}
\institute{Institution}
\date{\today}

% Remove the footer from title page
\setbeamertemplate{footline}{}


\title[Chapter Title]{Chapter: Introduction to Differential Equations}
\subtitle{Section 1.1 - Definitions and Terminology}
\author{Your Name}
\date{\today}

\begin{document}

% ============================================================================
% TITLE SLIDE
% ============================================================================
\begin{frame}
    \titlepage
\end{frame}

% ============================================================================
% OUTLINE SLIDE
% ============================================================================
\begin{frame}
    \frametitle{Outline}
    \tableofcontents
\end{frame}

% ============================================================================
% SECTION 1: DEFINITIONS
% ============================================================================
\section{Definitions}

\begin{frame}
    \frametitle{Section 1.1: Definitions and Terminology}
    
    In this section, we introduce fundamental concepts:
    \begin{itemize}
        \item Definition of differential equations
        \item Classification of differential equations
        \item Order and degree
    \end{itemize}
\end{frame}

\begin{frame}
    \frametitle{Definition 1: Differential Equation}
    
    \begin{definitionbox}{Differential Equation}
        A \textbf{differential equation} is an equation that contains one or more derivatives of an unknown function with respect to one or more independent variables.
        
        \vspace{0.3cm}
        
        General form: $F(x, y, y', y'', \ldots, y^{(n)}) = 0$
    \end{definitionbox}
    
    \pause
    
    \begin{examplebox}{Example}
        The equations:
        \begin{align*}
            \dydx &= 2x \\
            \dd2ydx2 + 5\dydx + 6y &= e^x
        \end{align*}
        are differential equations.
    \end{examplebox}
\end{frame}

\begin{frame}
    \frametitle{Definition 2: Order of a Differential Equation}
    
    \begin{definitionbox}{Order}
        The \textbf{order} of a differential equation is the order of the highest derivative present in the equation.
    \end{definitionbox}
    
    \pause
    
    \begin{examplebox}{Examples of Order}
        \begin{align*}
            \dydx &= 3x^2 \quad \text{(1st order)} \\
            \dd2ydx2 - 4y &= 0 \quad \text{(2nd order)} \\
            \frac{d^3y}{dx^3} + y &= \sin x \quad \text{(3rd order)}
        \end{align*}
    \end{examplebox}
\end{frame}

% ============================================================================
% SECTION 2: THEOREMS
% ============================================================================
\section{Theorems}

\begin{frame}
    \frametitle{Theorem 1: Existence and Uniqueness}
    
    \begin{theorembox}{Existence and Uniqueness Theorem}
        Let $R$ be a rectangular region in the $xy$-plane defined by $|x - x_0| \leq a$, $|y - y_0| \leq b$, and let
        \[
        f(x, y) \quad \text{and} \quad \frac{\partial f}{\partial y}
        \]
        be continuous on $R$. Then, there exists an interval $|x - x_0| \leq h$ (where $h > 0$) and a unique function $y(x)$ defined on this interval that satisfies the initial value problem:
        \[
        \dydx = f(x, y), \quad y(x_0) = y_0
        \]
    \end{theorembox}
\end{frame}

\begin{frame}
    \frametitle{Theorem 1 (continued)}
    
    \begin{examplebox}{Application}
        For the IVP:
        \begin{align*}
            \dydx &= \sqrt{y} \\
            y(0) &= 1
        \end{align*}
        
        We have $f(x,y) = \sqrt{y}$ and $\frac{\partial f}{\partial y} = \frac{1}{2\sqrt{y}}$.
        
        Both are continuous in a neighborhood of $(0, 1)$, so a unique solution exists near $x = 0$.
    \end{examplebox}
\end{frame}

% ============================================================================
% SECTION 3: WORKED EXAMPLES
% ============================================================================
\section{Worked Examples}

\begin{frame}
    \frametitle{Example 1: Verifying a Solution}
    
    \begin{problembox}{Problem}
        Verify that $y = e^{2x}$ is a solution to the differential equation:
        \[
        \dydx - 2y = 0
        \]
    \end{problembox}
\end{frame}

\begin{frame}
    \frametitle{Example 1: Solution}
    
    \begin{examplebox}{Solution}
        Given $y = e^{2x}$, we compute:
        \[
        \dydx = 2e^{2x}
        \]
        
        Substituting into the differential equation:
        \[
        \dydx - 2y = 2e^{2x} - 2e^{2x} = 0 \quad \checkmark
        \]
        
        Therefore, $y = e^{2x}$ is indeed a solution.
    \end{examplebox}
\end{frame}

\begin{frame}
    \frametitle{Example 2: Solving a Simple ODE}
    
    \begin{problembox}{Problem}
        Solve the differential equation:
        \[
        \dydx = 3x^2
        \]
    \end{problembox}
\end{frame}

\begin{frame}
    \frametitle{Example 2: Solution}
    
    \begin{examplebox}{Solution}
        Integrating both sides with respect to $x$:
        \begin{align*}
            \dydx &= 3x^2 \\
            \int dy &= \int 3x^2 \, dx \\
            y &= x^3 + C
        \end{align*}
        
        where $C$ is an arbitrary constant.
        
        \vspace{0.3cm}
        
        \textbf{General solution:} $y = x^3 + C$
    \end{examplebox}
\end{frame}

% ============================================================================
% CONCLUSION SLIDE
% ============================================================================
\section{Summary}

\begin{frame}
    \frametitle{Summary}
    
    \begin{itemize}
        \item A differential equation relates a function to its derivatives
        \item The order is determined by the highest derivative present
        \item The Existence and Uniqueness Theorem guarantees solutions under certain conditions
        \item Solutions can be verified by substitution
        \item Simple ODEs can often be solved by integration
    \end{itemize}
    
    \vspace{1cm}
    
    \centering
    \Large \textbf{Thank you!}
\end{frame}

\end{document}
